% -*- mode: latex; TeX-master: t; -*-

% \documentclass[notes=hide]{beamer}
% \documentclass[notes=only]{beamer}
\documentclass[
14pt,
aspectratio=169,
usenames,
dvipsnames,
% handout,
x11names]{beamer}
%\let\Tiny=\tiny

% Font selection
\overfullrule=5pt

\usepackage{etex}
\usepackage{pgfpages}

\usepackage{tabularx}
\usepackage{multicol}
% \usepackage[subfolder]{gnuplottex}

%% beamerscape dependencies
\usepackage[absolute,overlay]{textpos}
\setlength{\TPHorizModule}{\paperwidth}
\setlength{\TPVertModule}{\paperheight}
\textblockorigin{0mm}{0mm}

\definecolor{greenvm}{HTML}{AFD18E}
\definecolor{g}{HTML}{0E3B12}
\definecolor{v}{HTML}{003366}

\newcommand{\src}[1]{\scriptsize Figure Source: \textit{#1}}

%%%%%%%%%%%%%%%%%%%%%%%%%%%%%%%%%%%%%%%%%%%%%%%%%%%%%%%%%%%%%%%%%%%%%%%%%%%%%%%%
% Beamer template
\usetheme[
	bullet=square,             % Default option: circle
%	nobigpagenumber,         % No big circled page number on lower right
	noheadline,		 % No headline
%       nofootline,		 % No footline
	% grey,                    % For a neutral grey background color
%	watermark=BG_both,	 % png file for the watermark
]{Ithaca}
% \usetheme{metropolis}
% \usetheme{Boadilla}
%\usetheme{Berlin}
% \usetheme{Warsaw}
% % \usetheme{Dresden}
% %\usecolortheme{seagull}
% \usecolortheme{beaver}
% %\usecolortheme{seahorse}
%\usefonttheme{serif}
% \usefonttheme{structurebold}

%%%%%%%%%%%%%%%%%%%%%%%%%%%%%%%%%%%%%%%%%%%%%%%%%%%%%%%%%%%%%%%%%%%%%%%%%%%%%%%%
% Template
%%%%%%%%%%%%%%%%%%%%%%%%%%%%%%%%%%%%%%%%%%%%%%%%%%%%%%%%%%%%%%%%%%%%%%%%%%%%%%%%

\definecolor{mypurple}{HTML}{494949}
\definecolor{myyellow}{HTML}{FFCC33}
\definecolor{mygrey}{HTML}{999999}

\setbeamercolor{alerted text}{fg=myyellow}
\setbeamercolor{headline}{bg=footline.bg}
\setbeamercolor{block body alerted}{bg=normal text.bg!90!black}
\setbeamercolor{block body}{bg=normal text.bg!90!black}
\setbeamercolor{block body example}{bg=normal text.bg!90!black}
\setbeamercolor{block title alerted}{fg=myyellow}
\setbeamercolor{block title}{bg=blue}
\setbeamercolor{block title example}{use={normal text,example text},fg=example text.fg!75!normal text.fg,bg=normal text.bg!75!black}
\setbeamercolor{frametitle}{fg=white}
\setbeamercolor{item projected}{fg=white}
\setbeamercolor{normal text}{bg=mypurple,fg=white}
\setbeamercolor{section in sidebar}{fg=brown}
\setbeamercolor{section in sidebar shaded}{fg=grey}
\setbeamercolor{separation line}{}
\setbeamercolor{sidebar}{bg=red}
\setbeamercolor{sidebar}{parent=palette primary}
\setbeamercolor{structure}{bg=black, fg=white}
\setbeamercolor{subsection in sidebar}{fg=brown}
\setbeamercolor{subsection in sidebar shaded}{fg=grey}

% \setbeamercolor{palette primary}{fg=mypurple!80!black}
% \setbeamercolor{palette secondary}{fg=mypurple!80!black}
% \setbeamercolor{palette quaternary}{fg=mypurple!50!black}
% \setbeamercolor{page number in foot}{fg=white}

% lets add page numbers
% \expandafter\def\expandafter\insertshorttitle\expandafter{%
%   \insertshorttitle\hfill\insertframenumber\,/\,\inserttotalframenumber}

% \setbeamercovered{transparent}
\beamertemplatetransparentcoveredhigh

% *****************************************
% >>>> 15.3.2    Using Beamer's Templates
% *****************************************
% As a user of the beamer class you typically do not 'use' or 'invoke'
% templates yourself, directly. For example, the frame title template
% is automatically invoked by beamer somewhere deep inside the frame
% typesetting process. The same is true of most other
% templates. However, if, for whatever reason, you wish to invoke a
% template yourself, you can use the following command.
% \usebeamertemplate***{ element name }
% -------------------------------
%%% 7.2.1 The Headline and Footline
% \setbeamertemplate{headline} % Beamer-Template/-Color/-Font
% \setbeamertemplate{headline}
% {%
%   \begin{beamercolorbox}{section in head/foot}
%     \vskip2pt\insertnavigation{\paperwidth}\vskip2pt
%   \end{beamercolorbox}%
% }
% \setbeamertemplate{headline}[default] % The default is just an empty headline.
% \setbeamertemplate{headline}[infolines theme]
% \setbeamertemplate{headline}[miniframes theme]
% \setbeamertemplate{headline}[sidebar theme]
% \setbeamertemplate{headline}[smoothtree theme]
% \setbeamertemplate{headline}[smoothbars theme]
% \setbeamertemplate{headline}[tree]
% \setbeamertemplate{headline}[split theme]
% \setbeamertemplate{headline}[text line]{ text } % The headline is typeset with 'text'
% Logo in every slide except title
% \addtobeamertemplate{headline}{}
% {%
% \vspace{1ex}
% \hspace{1ex}
% \includegraphics[height=1cm]{myrmigki_color}
% }
% -------------------------------
% \setbeamertemplate{footline} % Beamer-Template/-Color/-Font
% \setbeamertemplate{footline}[default]
% \setbeamertemplate{footline}[infolines theme]
% \setbeamertemplate{footline}[miniframes theme]
% \setbeamertemplate{footline}[page number]
% \setbeamertemplate{footline}[frame number]
% \setbeamertemplate{footline}[split]
% \setbeamertemplate{footline}[text line]{ text }
% % Footline in every slideexcept title
\setbeamertemplate{footline}{
  \leavevmode%
  \begin{beamercolorbox}[left,wd=.4\paperwidth,ht=2.5ex,dp=1.125ex,leftskip=.3cm, rightskip=.3cm plus1fil]{footline}%
    \insertshorttitle%
  \end{beamercolorbox}
  %
  \begin{beamercolorbox}[center,wd=.2\paperwidth,ht=2.5ex,dp=1.125ex,leftskip=.3cm,rightskip=.3cm plus1fil]{footline}%
    \centering
    \insertframenumber~of~\inserttotalframenumber%
  \end{beamercolorbox}%
  %
  \begin{beamercolorbox}[right,wd=.4\paperwidth,ht=2.5ex,dp=1.125ex,leftskip=.3cm,rightskip=.3cm plus1fil]{footline}%
    \hskip2ex plus1fill%
    \includegraphics[height=1.5ex]{cc}~%
    \includegraphics[height=1.5ex]{by}~~%
    % \includegraphics[height=1.5ex]{sa}~~%
    \insertshortauthor~%
  \end{beamercolorbox}%
}%
% -------------------------------
%%% 7.2.2 The Sidebars
% -------------------------------
%%% 7.2.3 Navigation Bars (funktioniert nur mit miniframe Themes)
% \setbeamertemplate{mini frames}[default] % shows small circles as mini frames.
% \setbeamertemplate{mini frames}[box] % shows small rectangles as mini frames.
% \setbeamertemplate{mini frames}[tick] % shows small vertical bars as mini frames.
% -------------------------------
%%% 7.2.4 The Navigation Symbols
%%% Beamer-Template/-Color/-Font navigation symbols
% \setbeamertemplate{navigation symbols}{} % suppresses all navigation symbols:
% \setbeamertemplate{navigation symbols}[horizontal] % Organizes the navigation symbols horizontally.
% \setbeamertemplate{navigation symbols}[vertical] % Organizes the navigation symbols vertically.
% \setbeamertemplate{navigation symbols}[only frame symbol] % Shows only the navigational symbol for navigating frames.
% -------------------------------
%%% 7.2.5 The Logo
% \setbeamertemplate{logo} % Beamer-Template/-Color/-Font
% -------------------------------
%%% 7.2.6 The Frame Title
% \setbeamertemplate{frametitle} % Beamer-Template/-Color/-Font
% \setbeamertemplate{frametitle}[default][center] % left, center, right
% \setbeamertemplate{frametitle}[shadow theme]
% \setbeamertemplate{frametitle}[sidebar theme]
% \setbeamertemplate{frametitle}[smoothbars theme]
% \setbeamertemplate{frametitle}[smoothtree theme]
% \setbeamertemplate{frametitle}
% {
%     \nointerlineskip
%     \begin{beamercolorbox}[sep=0.3cm,ht=2.5em,wd=\paperwidth]{frametitle}
%         \vbox{}\vskip-2ex%
%         \hspace{1.2cm}
%         \strut\insertframetitle\strut
%         \vskip-0.8ex%
%     \end{beamercolorbox}
% }
% % Minimalistic frametitle
% \setbeamertemplate{frametitle}
% {
%   \vspace{.2cm}
%   \hspace{-.8cm}
%   \insertframetitle
% }
% -------------------------------
%%% 7.2.7 The Background
% \setbeamertemplate{background canvas} % Beamer-Template/-Color/-Font
% \setbeamertemplate{background canvas}[default]
% \setbeamertemplate{background canvas}[vertical shading][ color options ] installs a vertically shaded background.
%     - top= color specifies the color at the top of the page. By default, 25% of the foreground of
%       the beamer-color palette primary is used.
%     - bottom= color specifies the color at the bottom of the page. By default, the background of
%       normal text at the moment of invocation of this command is used.
%     - middle= color specifies the color for the middle of the page. Thus, if this option is given, the
%       shading changes from the bottom color to this color and then to the top color.
%     - midpoint= factor specifies at which point of the page the middle color is used. A factor of 0
%       is the bottom of the page, a factor of 1 is the top. The default, which is 0.5 is in the middle.
% \setbeamertemplate{background canvas}[vertical shading][top=gray!40,bottom=gray!40,middle=white,midpoint=.9]
%\setbeamertemplate{background canvas}[vertical shading][top=white,bottom=gray!40]
% \setbeamercolor{background canvas}{bg=white}
%% Slide Background
\setbeamertemplate{background canvas}{%
	\begin{tikzpicture}
		\path [outer color = mypurple,
		       inner color = mypurple!75!white]
			(0,0) rectangle (\paperwidth,\paperheight);
	\end{tikzpicture}%
}

% \setbeamertemplate{background} % Beamer-Template/-Color/-Font
% \setbeamertemplate{background}[default] % is empty.
% \setbeamertemplate{background}[grid][step=1cm] % places a grid on the background.
%     - step= dimension specifies the distance between grid lines. The default is 0.5cm.
%     - color= color specifies the color of the grid lines. The default is 10% foreground.
% % Add global transition effect \transfade (breaks the background)
% \addtobeamertemplate{background canvas}{%
%   \transfade[duration=.1]%
% }{}
% -------------------------------
%%% 7.3 Margin Sizes
\setbeamersize{text margin left=2em,text margin right=2em}
% \setbeamersize{sidebar width left=2cm}
%         - text margin left= TEX dimension sets a new left margin. This excludes the left sidebar. Thus,
%           it is the distance between the right edge of the left sidebar and the left edge of the text.
%         - text margin right= TEX dimension sets a new right margin.
%         - sidebar width left= TEX dimension sets the size of the left sidebar. Currently, this command
%           should be given before a shading is installed for the sidebar canvas.
%         - sidebar width right= TEX dimension sets the size of the right sidebar.
%         - description width= TEX dimension sets the default width of description labels, see Section 11.1.
%         - description width of= text sets the default width of description labels to the width of the
%             text , see Section 11.1.
%         - mini frame size= TEX dimension sets the size of mini frames in a navigation bar. When two
%           mini frame icons are shown alongside each other, their left end points are TEX dimension far
%           apart.
%         - mini frame offset= TEX dimension set an additional vertical offset that is added to the mini
%           frame size when arranging mini frames vertically.
% -------------------------------
%%% 9.1 Adding a Title Page
% \setbeamersize{title page} % Beamer-Template/-Color/-Font
%    This template is invoked when the \titlepage command is used.
%    The following commands are useful for this template:
%     -  \insertauthor inserts a version of the author's name that is useful for the title page.
%     -  \insertdate inserts the date.
%     -  \insertinstitute inserts the institute.
%     -  \inserttitle inserts a version of the document title that is useful for the title page.
%     -  \insertsubtitle inserts a version of the document title that is useful for the title page.
%     -  \inserttitlegraphic inserts the title graphic into a template.
\defbeamertemplate*{title page}{flat2}[1][]
{
  \addtocounter{framenumber}{-1}
  \begin{center}
    \Huge{\inserttitle}\\[1ex]
    \hrule
    \vspace{1ex}
    \Large{\insertsubtitle}
    \vfill

    % \large\textcolor{normal text.fg}{\textit{\insertinstitute}, \insertdate}
    \vspace{1em}

    \vfill

    \large \insertauthor\\
    \vspace{-1ex}
    {\scriptsize \insertdate}\\
    \insertlogo
  \end{center}
}

% -------------------------------
%%% 9.2 Adding Sections and Subsections
% -------------------------------
%%% Parent Beamer-Template sections/subsections in toc
% This is a parent template, whose children are section in toc and subsection in toc.
% \setbeamertemplate{sections/subsections in toc}[default]
% \setbeamertemplate{sections/subsections in toc}[sections numbered]
% \setbeamertemplate{sections/subsections in toc}[subsections numbered]
% \setbeamertemplate{sections/subsections in toc}[circle]
% \setbeamertemplate{sections/subsections in toc}[square]
% \setbeamertemplate{sections/subsections in toc}[ball]
% \setbeamertemplate{sections/subsections in toc}[ball unnumbered]
% -------------------------------
%%% 9.6 Adding a Bibliography
% -------------------------------
% \setbeamertemplate{bibliography item} % Beamer-Template/-Color/-Font
% \setbeamertemplate{bibliography item}[default] %  little article icon as the reference
% \setbeamertemplate{bibliography item}[article] % Alias for the default.
% \setbeamertemplate{bibliography item}[book] % little book icon as the reference
% \setbeamertemplate{bibliography item}[triangle] % triangle as the reference
% \setbeamertemplate{bibliography item}[text] % reference text (like '[Dijkstra, 1982]')
% -------------------------------
%%% 10.1 Adding Hyperlinks and Buttons
% -------------------------------
%%% 11.1 Itemizations, Enumerations, and Descriptions
% \setbeamertemplate{items} % parent template of itemize items and enumerate items
% \setbeamertemplate{itemize items} % Parent Beamer-Template
% \setbeamertemplate{itemize items}[triangle]
% \setbeamertemplate{itemize items}[circle]
% \setbeamertemplate{itemize items}[square]
% \setbeamertemplate{itemize items}[ball]
% -------------------------------
% \setbeamertemplate{enumerate items}[default] % Numbered
% \setbeamertemplate{enumerate items}[circle] % Places the numbers inside little circles.
% \setbeamertemplate{enumerate items}[square] % Places the numbers on little squares.
% \setbeamertemplate{enumerate items}[ball] % 'Projects' the numbers onto little balls.
% -------------------------------
%%% 11.2 Hilighting
% -------------------------------
%%% 11.3 Block Environments
% \setbeamertemplate{blocks} % Parent Beamer-Template
% \setbeamertemplate{blocks}[default]
% \setbeamertemplate{blocks}[rounded][shadow=true]
 % \setbeamertemplate{blocks}[rounded][shadow=false]
% -------------------------------
%%% 11.4 Theorem Environments
% \setbeamertemplate{qed symbol} % Beamer-Template/-Color/-Font
% -------------------------------
% \setbeamertemplate{theorems} % Parent Beamer-Template
% \setbeamertemplate{theorems}[default]
% \setbeamertemplate{theorems}[normal font]
% \setbeamertemplate{theorems}[numbered]
% \setbeamertemplate{theorems}[ams style]
% -------------------------------
%%% 11.6 Figures and Tables
% \setbeamertemplate{caption} % Beamer-Template/-Color/-Font
% \setbeamertemplate{caption}[default] typesets the caption name (a word like 'Figure' or 'Abbildung' or 'Table')
% \setbeamertemplate{caption}[numbered] adds the figure or table number to the caption.
% \setbeamertemplate{caption}
% -------------------------------
% \setbeamertemplate{caption name} % Beamer-Color/-Font
% -------------------------------
%%% 11.10    Abstract
% -------------------------------
%%% 11.11 Verse, Quotations, Quotes
% -------------------------------
%%% 11.12 Footnotes
% -------------------------------
%%% 18.1 Specifying Note Contents
% \setbeamertemplate{note page} % Beamer-Template/-Color/-Font
% \setbeamertemplate{note page}[default]
% \setbeamertemplate{note page}[compress]
% \setbeamertemplate{note page}[plain]
% -------------------------------
%%% Specifying Which Notes and Frames Are Shown
 % \setbeameroption{hide notes}
% \setbeameroption{show notes}
% \setbeameroption{show notes on second screen= right }
% \setbeameroption{show only notes}


\renewcommand{\footnoterule}{}

\usepackage{hyperref}
\hypersetup{
  linktoc=all,
  bookmarks=false,           % show bookmarks bar?
  bookmarksopen,
  bookmarksnumbered,
  colorlinks = false,
  linkcolor=black,           % color of interlinks
  citecolor=black,           % color of the sitation links
  urlcolor=black,            % the url color
  unicode=true,              % non-Latin characters in Acrobat’s bookmarks
  pdftoolbar=true,           % show Acrobat’s toolbar?
  pdfmenubar=true,           % show Acrobat’s menu?
  pdffitwindow=true,         % window fit to page when opened
  pdftitle={Managed Runtime Systems},  % title
  pdfborder={ 0 0 0 },        % uBorder tin links
  pdfauthor = {Foivos Zakkak},
  pdfcreator = {Foivos Zakkak},
}

\usefonttheme{professionalfonts}% use own font handling
\usepackage{fontspec}
\usepackage{xunicode}
% \setmainfont{Comfortaa}
% \setsansfont{Comfortaa}
% \setsansfont{Liberation Sans}
\setmainfont[
% SmallCapsFont={Linux Biolinum},
SmallCapsFeatures={Letters=SmallCaps},
]{Liberation Sans}
% \setmonofont{Liberation Mono}
% \setmonofont{Inconsolata LGC for Powerline}
% Math fonts
\usepackage{unicode-math}
\setmathfont{xits-math.otf}

\usepackage{graphicx}
\usepackage{ulem}
\usepackage{color}
\usepackage{xspace}
\usepackage{xcolor}
\usepackage{array}
\usepackage{tikz}
\usetikzlibrary{arrows,arrows.meta,shapes,decorations,decorations.pathmorphing,calc,shadows,shadows.blur,shapes.multipart,positioning,patterns,fit,backgrounds,trees,shapes}

\tikzset{
    %Define standard arrow tip
    >=stealth',
    % Define arrow style
    arrow/.style={
           ->,
           thick,
           color=uompurple,
           shorten <=2pt,
           shorten >=4pt,},
    % Daniel's proposal for "uncovering" parts of a tikz-tree %
    invisible/.style={opacity=0},
    visible on/.style={alt=#1{}{invisible}},
    alt/.code args={<#1>#2#3}{%
      \alt<#1>{\pgfkeysalso{#2}}{\pgfkeysalso{#3}} % \pgfkeysalso doesn't change the path
    },
  }
\pgfdeclaredecoration{penciline}{initial}{
    \state{initial}[width=+\pgfdecoratedinputsegmentremainingdistance,
    auto corner on length=1mm,]{
        \pgfpathcurveto%
        {% From
            \pgfqpoint{\pgfdecoratedinputsegmentremainingdistance}
                      {\pgfdecorationsegmentamplitude}
        }
        {%  Control 1
        \pgfmathrand
        \pgfpointadd{\pgfqpoint{\pgfdecoratedinputsegmentremainingdistance}{0pt}}
                    {\pgfqpoint{-\pgfdecorationsegmentaspect
                     \pgfdecoratedinputsegmentremainingdistance}%
                               {\pgfmathresult\pgfdecorationsegmentamplitude}
                    }
        }
        {%TO
        \pgfpointadd{\pgfpointdecoratedinputsegmentlast}{\pgfpoint{1pt}{1pt}}
        }
    }
    \state{final}{}
  }
\usepackage{amsmath}
\usepackage{textpos}
% \usepackage{subfigure}

\graphicspath{{figures/}}
% \newcommand{\longname}{\emph{Source-Level Compiler Optimizations for Task-Parallelism}}
\newcommand{\java}[0]{Java\texttrademark\xspace}
\newcommand{\code}[1]{\texttt{#1}}

\usepackage{bbding} % for the Checkmark and XSolidBrush
\newcommand{\tik}[0]{{\color{green}\Checkmark}} % Check-mark
\newcommand{\ex}[0]{{\color{red}\XSolidBrush}}  % X-mark

% \usepackage{verbatim}
\usepackage{listings}
\lstset{
  language=c,
  basicstyle=\scriptsize\ttfamily,
%  columns=flexible,
  % numbers=left,
  % numberstyle=\ttfamily\tiny,
  showstringspaces=false,
  alsoletter={-},
  literate={-}{-}1,
%  numbersep=1em,
  xleftmargin=2em,
%  xrightmargin=1em,
  escapeinside={@}{@},
%  morecomment=[l][\color{BrickRed}]{\#pragma\ task},
  commentstyle=\color{black},
  morekeywords={final, transient, synthetic},
%   identifierstyle=\color{green},
%  backgroundcolor=\color{Honeydew1},
  keywordstyle=\bfseries\color{myyellow},
  stringstyle=\color{YellowGreen},
%  moredelim=**[is][\only<+(1)->{\color{black}\lstset{style=highlight}}]{~}{~},
}

\usepackage{booktabs} % bottomrule
\usepackage{multirow}

% Print table of contects at the beggining of each section
% \AtBeginSection[]{\begin{frame}\frametitle{Table of Contents}\tableofcontents[currentsection,currentsubsection]\end{frame}}
% \AtBeginSubsection[]{\begin{frame}\frametitle{Table of Contents}\tableofcontents[currentsection,currentsubsection]\end{frame}}

% Add outlines
% \AtBeginSection[]
% {
%   {
%   \setbeamertemplate{footline}{}
%   \begin{frame}<beamer>[plain,noframenumbering]
%     \frametitle{Outline}
%     \tableofcontents[currentsection]
%   \end{frame}
%   }
% }
\setcounter{tocdepth}{1}

%%%%%%%%%%%%%%%
% Intro slide %
%%%%%%%%%%%%%%%

\title{Managed Runtime Systems}
\subtitle{Lecture 02: The \java Virtual Machine}
\author[\url{https://foivos.zakkak.net}]{Foivos Zakkak}
% \institute[]{ManLang'17}
\date{\url{https://foivos.zakkak.net}}
% \logo{\includegraphics[height=1.0cm]{uniman-logo}}
%% Spread those bullets
% \let\olditem\item
% \renewcommand{\item}{\setlength{\itemsep}{\fill}\olditem}

\begin{document}
\setbeamercovered{invisible}

% \maketitle

\begin{frame}[plain]
  \titlepage
  \centering
  \includegraphics[height=.75cm]{cc}~
  \includegraphics[height=.75cm]{by}\\[1em]
  % \includegraphics[height=.75cm]{nc-eu}~
  % \includegraphics[height=.75cm]{sa}\\[1em]
  \scriptsize{Except where otherwise noted, this presentation is licensed under the\\
    \href{http://creativecommons.org/licenses/by/4.0/}%
    {Creative Commons Attribution 4.0 International License.}\\[1ex]
    Third party marks and brands are the property of their respective
    holders.}
\end{frame}

% \begin{frame}[plain,noframenumbering]{Outline}
%     \tableofcontents
% \end{frame}

\section{The Java Virtual Machine}

\begin{frame}[fragile]{The JVM Architecture}
  \centering

  \tikzstyle{file} = [decorate,text=black, fill=Dandelion, rectangle, minimum width=1cm, minimum height=1.7cm]
  \tikzstyle{arrow} = [->,thick]
  \tikzstyle{darrow} = [<->,thick]
  \tikzstyle{component} = [text=black, fill=white, rectangle, minimum width=2cm, minimum height=1cm]
  \tikzstyle{memory} = [component]

  \scriptsize
  \begin{tikzpicture}[decoration=penciline]
    \node (source) [file] {.c};
    \node (bytecode) [file, right=of source] {.class};

    \draw[arrow] (source) -- (bytecode) node[midway,above] {\texttt{javac}};

    \node (classloader) [component, right=of bytecode, fill=BrickRed, text=white] {Class Loader};

    \draw[arrow] (bytecode) -- (classloader) node[midway,above] {\texttt{java}};

    \node (methoda) [memory, below left=of bytecode] {Method Area};
    \node (heapa) [memory, right=5mm of methoda] {Heap Area};
    \node (stacka) [memory, right=5mm of heapa] {Stack Area};
    \node (pcregistera) [memory, right=5mm of stacka] {PC Register};
    \node (nativea) [memory, right=5mm of pcregistera] {Native Stack};

    \node[fit=(methoda)(nativea), draw=YellowGreen, thick, inner sep=3mm] (data) {};
    \node[text=YellowGreen, thick, above=3mm of methoda] {\small Data locations};

    \draw[darrow] (classloader) -- (data);

    \node (exec) [memory, below=of heapa, fill=BrickRed, text=white] {Execution Engine};
    \node (nativem) [memory, right=of exec] {Native Methods Interface};
    \node (nativelib) [memory, right=of nativem] {Native Libraries};

    \draw[darrow] (exec) -- ($ (exec.north) + (0,.7) $);
    \draw[darrow] (exec) -- (nativem);
    \draw[darrow] (nativem) -- ($ (nativem.north) + (0,.7) $);
    \draw[arrow] (nativelib) -- (nativem);

  \end{tikzpicture}
\end{frame}

\begin{frame}{Class loading}
  \begin{enumerate}
  \item Verification
  \item Preparation of JVM internal data structures
  \item Resolution of references (might load other classes)
  \item Initialization
  \end{enumerate}
\end{frame}

\note[itemize]{
\item Loading: A class is obtained from a file and loaded to the JVM memory.
\item Verifying: Check whether or not the read class is configured as described in the Java Language Specification and JVM specifications. This is the most complicated test process of the class load processes, and takes the longest time. Most cases of the JVM TCK test cases are to test whether or not a verification error occurs by loading wrong classes.
\item Preparing: Prepare a data structure that assigns the memory required by classes and indicates the fields, methods, and interfaces defined in the class.
\item Resolving: Change all symbolic references in the constant pool of the class to direct references.
\item Initializing: Initialize the class variables to proper values. Execute the static initializers and initialize the static fields to the configured values.
}

\begin{frame}[fragile]{Java Class file}
  \begin{columns}
    \begin{column}{.7\linewidth}
\begin{lstlisting}
ClassFile {
  u4             magic;
  u2             minor_version;
  u2             major_version;
  u2             constant_pool_count;
  cp_info        constant_pool[constant_pool_count-1];
  u2             access_flags;
  u2             this_class;
  u2             super_class;
  u2             interfaces_count;
  u2             interfaces[interfaces_count];
  u2             fields_count;
  field_info     fields[fields_count];
  u2             methods_count;
  method_info    methods[methods_count];
  u2             attributes_count;
  attribute_info attributes[attributes_count];
}
\end{lstlisting}
    \end{column}
    \begin{column}{.3\linewidth}
      \centering

      \tikzstyle{stack} = [rectangle, draw, minimum width=3cm, thick]
      \tikzstyle{arrow} = [->, thick]

      \scriptsize
      \begin{tikzpicture}
        \node[stack] (f0) {magic};
        \node[stack, below=0mm of f0] (f1) {minor-major version};
        % \node[stack, below=0mm of f1] (f2) {major version};
        \node[stack, below=0mm of f1] (f3) {constant pool count};
        \node[stack, below=0mm of f3, draw=YellowGreen, text=YellowGreen] (f4) {constant pool[]};
        \node[stack, below=0mm of f4] (f5) {access flags};
        \node[stack, below=0mm of f5] (f6) {this class};
        \node[stack, below=0mm of f6] (f7) {super class};
        \node[stack, below=0mm of f7] (f8) {interfaces count};
        \node[stack, below=0mm of f8] (f9) {interfaces[]};
        \node[stack, below=0mm of f9] (f10) {fields count};
        \node[stack, below=0mm of f10] (f11) {fields[]};
        \node[stack, below=0mm of f11] (f12) {methods count};
        \node[stack, below=0mm of f12] (f13) {methods[]};
        \node[stack, below=0mm of f13] (f14) {attributes count};
        \node[stack, below=0mm of f14] (f15) {attributes[]};

        \draw[arrow] (f6.east) to[bend right] (f4.east);
        \draw[arrow] (f7.east) to[bend right] (f4.east);
        \draw[arrow] (f9.east) to[bend right] (f4.east);

      \end{tikzpicture}
    \end{column}
\end{columns}
\end{frame}

\begin{frame}[fragile]{Field Info}
\begin{lstlisting}
field_info {
  u2 access_flags; // public, private, protected, static, final,
                   // volatile, transient, synthetic, enum
  u2 name_index;
  u2 descriptor_index;
  u2 attributes_count;
  attribute_info attributes[attributes_count]; // e.g., Deprecated
}
\end{lstlisting}
\end{frame}

\begin{frame}{Java fields}
  \begin{itemize}  \setlength{\itemsep}{\fill}
  \item Reference (Size is implementation-dependent)
  \item Primitive
    \begin{itemize}
    \item 8, 16, 32, or 64-bit Integer (\texttt{byte}, \texttt{short}, \texttt{int}, \texttt{long}, \texttt{boolean}, \texttt{char})
    \item 32 or 64-bit IEEE 754 floating-point (\texttt{float}, \texttt{double})
    \end{itemize}
  \end{itemize}
\end{frame}

\begin{frame}{Field Descriptors}
  \begin{tabular}{l | l | l}
    B                      & byte      &                                                               \\
    C                      & char      &                                                               \\
    D                      & double    &                                                               \\
    F                      & float     &                                                               \\
    I                      & int       &                                                               \\
    J                      & long      &                                                               \\
    L \textit{ClassName} ; & reference & \texttt{Ljava/lang/String;}                                   \\
    S                      & short     &                                                               \\
    Z                      & boolean   &                                                               \\\relax
    [                      & array     & \texttt{[I} ~~ \texttt{[[[D} ~~  \texttt{[Ljava/lang/Object;} \\
  \end{tabular}
\end{frame}

\begin{frame}[fragile]{Method Info}
\begin{lstlisting}
method_info {
  u2 access_flags; // public, private, native, var args, synthetic ...
  u2 name_index;
  u2 descriptor_index;
  u2 attributes_count;
  attribute_info attributes[attributes_count]; // e.g., Code, Exceptions
}
\end{lstlisting}
\begin{exampleblock}{Method Descriptors}
  \texttt{void foo(int a, Object o)} => \texttt{(ILjava/lang/Object;)V}\\
  \texttt{String bar(int[][] i)} => \texttt{([[I)Ljava/lang/String;}
\end{exampleblock}
\end{frame}

\begin{frame}[fragile]{Code attribute}
\begin{lstlisting}
Code_attribute {
  u2 attribute_name_index;
  u4 attribute_length;
  u2 max_stack;
  u2 max_locals;
  u4 code_length;
  u1 code[code_length]; // The actual bytecode
  u2 exception_table_length;
  { u2 start_pc; // if catch_type is thrown between start_pc and end_pc,
    u2 end_pc;   // execute handler at handler_pc
    u2 handler_pc;
    u2 catch_type;
  } exception_table[exception_table_length];
  u2 attributes_count;
  attribute_info attributes[attributes_count]; // line numbers, stack maps,...
}
\end{lstlisting}
\end{frame}

\section{Runtime}

\begin{frame}{Java Stack vs Native Stack}
  \centering

  \tikzstyle{frame} = [rectangle, draw, thick, minimum width=2cm, minimum height=.75cm]
  \tikzstyle{stack} = [rectangle, draw, minimum width=2.5cm]

  \scriptsize
  \begin{tikzpicture}
    \node (title) {Java Stack};

    \node[frame, below=of title] (frame1) {frame};
    \node[frame, below=2mm of frame1, draw=YellowGreen] (frame2) {frame};
    \node[frame, below=2mm of frame2] (frame3) {frame};

    \node[fit=(title)(frame3), draw, thick, inner sep=3mm] (jstack) {};

    \draw[->, thick] ($(frame1.north) + (0,2mm) $) -- ($(title.south) - (0,2mm) $);

    \node[right=2cm of jstack.north east, below, align=center] (frametitle) {~\\[2ex]Frame};

    \node[stack, below=of frametitle] (oper2) {operand stack 2};
    \node[stack, below=0mm of oper2] (oper1) {operand stack 1};
    \node[stack, below=0mm of oper1] (oper0) {operand stack 0};
    \node[stack, below=0mm of oper0] (localn) {local variable n};
    \node[stack, below=0mm of localn] (localdot) {\large$\vdots$};
    \node[stack, below=0mm of localdot] (local1) {local variable 1};
    \node[stack, below=0mm of local1] (local0) {local variable 0};
    \node[stack, below=0mm of local0] (ra) {return address};

    \node[fit=(frametitle)(ra), draw=YellowGreen, thick, inner sep=0mm] (jstack) {};

    \draw[->, thick] ($(oper2.north) + (0,2mm) $) -- ($(frametitle.south) - (0,2mm) $);

    \draw[thick, dashed, draw=YellowGreen] (frame2.north east) -- (jstack.north west);
    \draw[thick, dashed, draw=YellowGreen] (frame2.south east) -- (jstack.south west);

  \end{tikzpicture}
\end{frame}

\begin{frame}{Java Threads}
  \begin{itemize}
  \item A Program Counter\footnote{Class+method+bytecode-offset or native code pointer}
  \item A Java Stack
  \item Native stack
  \end{itemize}
\end{frame}

\begin{frame}{Java Objects and Arrays}
  \centering

  \tikzstyle{field} = [rectangle, draw, thick, minimum width=2.5cm]

  \scriptsize
  \begin{tikzpicture}
    \node (object) {\textbf{Object}};

    \node[field, below=3mm of object, fill=PineGreen] (cheader) {header};
    \node[field, below=0mm of cheader] (field0) {field 0};
    \node[field, below=0mm of field0] (field1) {field 1};
    \node[field, below=0mm of field1] (fields) {\large$\vdots$};
    \node[field, below=0mm of fields] (fieldn) {field n};

    \node[right=4cm of object] (array) {\textbf{Array}};

    \node[field, below=3mm of array, fill=PineGreen] (aheader) {header};
    \node[field, below=0mm of aheader] (length) {length};
    \node[field, below=0mm of length] (element0) {element 0};
    \node[field, below=0mm of element0] (element1) {element 1};
    \node[field, below=0mm of element1] (elements) {\large$\vdots$};
    \node[field, below=0mm of elements] (elementn) {element n};
  \end{tikzpicture}
\end{frame}


% Closing slide
% \subsection{Closure}
% % {
% % \setbeamertemplate{footline}{}
% % \begin{frame}[noframenumbering]
% %  \frametitle{\fontspec{Purisa}\textbf{Thank You!}}
% %  \centering
% %  \titlepage
% % \end{frame}
% % }
% \begin{frame}
%   % \centering
%   \Huge
%   \alert{Can managed programming languages run on future
%   many-core architectures?}
%   % \vspace{-2em}
%   \vfill
%   \begin{flushright}
%     \fontspec{Purisa}\structure{Thank You!}
%   \end{flushright}
% \end{frame}

%   BACKUP SLIDES

\end{document}
